\chapter{Analyse og kravspecifikation}
\label{chap:analysis}

\section{Domæneanalyse}
\label{sec:domæneanalyse}

For at forstå det problem systemet skal løse, er det nødvendigt at analysere den kontekst, hvori vedligeholdelsesteknikeren arbejder. Denne analyse tager udgangspunkt i information indsamlet i samarbejde med Arkyn Studios samt de erfaringer der er gjort i løbet af praktikopholdet.

\subsection{Teknikerens arbejdssituation}
\label{subsec:teknikerens-arbejdssituation}

En vedligeholdelsesteknikers arbejdsdag er præget af uforudsigelige hændelser: maskiner der går ned, fejl der skal diagnosticeres og udstyr der skal inspiceres. Arbejdet foregår på fabriksgulvet — typisk under tidspres, i støjende omgivelser og med begrænset adgang til stationære computere og dokumentation.

Når en teknikeren møder en ukendt maskine eller et ukendt udstyr, følger typisk dette forløb:

\begin{enumerate}
  \item \textbf{Identifikation:} Teknikeren forsøger at identificere maskinen ved at aflæse typeskilt, mærkat eller stregkode — hvis disse overhovedet er læselige efter års slitage og industrielt miljø.
  \item \textbf{Informationssøgning:} På baggrund af den identificerede model eller serienummer søges i virksomhedens dokumentationssystem, papirmanualer eller via kolleger.
  \item \textbf{Udførelse:} Teknikeren udfører vedligeholdelsen baseret på fundet information — eller på bedste skøn, hvis informationen er utilstrækkelig.
\end{enumerate}

Dette forløb er forbundet med kendte smertepunkter: typeskilte kan være ulæselige eller på fremmedsprog, dokumentation kan være svær at finde, og afhængighed af erfarne kolleger skaber flaskehalse i organisationen.

\subsection{Teknologisk kontekst}
\label{subsec:teknologisk-kontekst}

Centrale teknologier i løsningsrummet er:

\textbf{Optical Character Recognition (OCR):} Teknologi til automatisk tekstgenkendelse i billeder. I industriel kontekst anvendes OCR til aflæsning af typebetegnelser, serienumre og tekniske specifikationer fra maskinlabels. Moderne cloud-baserede OCR-løsninger som Google Cloud Vision~\cite{google_vision} opnår høj nøjagtighed selv på skæve vinkler og delvist slidte labels.

\textbf{Large Language Models (LLM):} AI-modeller der kan fortolke struktureret tekstinformation og generere handlingsorienterede vejledninger på naturligt sprog. Modeller som Anthropic Claude~\cite{anthropic_claude} kan instrueres til at returnere struktureret JSON-output, hvilket letter integration med eksisterende systemer.

\textbf{Mobilkamera:} Moderne smartphones har kameraer af tilstrækkelig kvalitet til at fungere som input-enhed for OCR-behandling, og giver teknikeren et allerede velkendt og tilgængeligt redskab.

Selvom de individuelle teknologier er velkendte, er kombinationen af billedbaseret input, realtids-OCR og AI-assisteret vejledning i en vedligeholdelseskontekst et relativt uudforsket område.

\section{Brugsscenarier}
\label{sec:brugsscenarier}

På baggrund af domæneanalysen er følgende primære brugsscenarier identificeret.

\subsection{Use Case 1: Billede af typeskilt}

\begin{tabular}{@{}lp{0.7\textwidth}@{}}
  \toprule
  \textbf{Aktør}        & Vedligeholdelsesteknikeren \\
  \textbf{Forudsætning} & Teknikeren er i nærheden af en maskine med et synligt typeskilt \\
  \midrule
  \multicolumn{2}{@{}l@{}}{\textbf{Forløb}} \\
  1 & Teknikeren åbner applikationen og navigerer til kamera-fanen \\
  2 & Der tages et billede af maskinens typeskilt \\
  3 & Systemet kører OCR og returnerer strukturerede entiteter (producent, model, serienummer mv.) \\
  4 & Teknikeren trykker \textit{Analyser} på den ønskede entitet \\
  5 & Systemet returnerer AI-genereret vedligeholdelsesvejledning \\
  \midrule
  \textbf{Resultat} & Teknikeren modtager relevant information inden for få sekunder \\
  \bottomrule
\end{tabular}

\vspace{1em}

\subsection{Use Case 2: Video af stor label}

\begin{tabular}{@{}lp{0.7\textwidth}@{}}
  \toprule
  \textbf{Aktør}        & Vedligeholdelsesteknikeren \\
  \textbf{Forudsætning} & Maskinens label er for stor til at blive fotograferet i ét billede \\
  \midrule
  \multicolumn{2}{@{}l@{}}{\textbf{Forløb}} \\
  1 & Teknikeren optager en kort video ved at panorere kameraet henover labelen \\
  2 & Systemet behandler videoen ved at udtrække frames (1 FPS) og køre OCR på hver enkelt \\
  3 & Systemet aggregerer den udtrukne tekst og returnerer strukturerede entiteter \\
  4--5 & Som Use Case 1, trin 4--5 \\
  \midrule
  \textbf{Resultat} & Komplet tekstudtræk fra labels der er for store til ét billede \\
  \bottomrule
\end{tabular}

\vspace{1em}

\subsection{Use Case 3: Lav billedkvalitet}

\begin{tabular}{@{}lp{0.7\textwidth}@{}}
  \toprule
  \textbf{Aktør}        & Vedligeholdelsesteknikeren \\
  \textbf{Forudsætning} & Billede er optaget under dårlige lysforhold, eller typeskiltet er snavset \\
  \midrule
  \multicolumn{2}{@{}l@{}}{\textbf{Forløb}} \\
  1 & Teknikeren uploader et billede som normalt \\
  2 & Systemet registrerer lav OCR-confidence (under 60~\%) \\
  3 & Systemet markerer resultatet som \textit{Manuel kontrol påkrævet} \\
  4 & Teknikeren informeres om at tage et nyt, tydeligere billede \\
  \midrule
  \textbf{Resultat} & Systemet fejler kontrolleret frem for at returnere forkerte data \\
  \bottomrule
\end{tabular}

\section{Kravspecifikation}
\label{sec:krav}

\subsection{Funktionelle krav}
\label{subsec:funktionelle-krav}

Tabel~\ref{tab:funktionelle-krav} opstiller de funktionelle krav til systemet, udledt fra domæneanalysen og brugsscenariet.

\begin{table}[H]
  \centering
  \caption{Funktionelle krav}
  \label{tab:funktionelle-krav}
  \begin{tabularx}{\textwidth}{@{}llX@{}}
    \toprule
    \textbf{ID} & \textbf{Krav} & \textbf{Beskrivelse} \\
    \midrule
    F1 & Kameraoptagelse     & Systemet skal understøtte optagelse af billeder og videoer via enhedens kamera \\
    \addlinespace
    F2 & OCR-udtræk          & Systemet skal udtrække tekst fra uploadede billeder og videoer via Google Cloud Vision \\
    \addlinespace
    F3 & Entitetsformatering & Systemet skal formatere råtekst til strukturerede entiteter med tilhørende handlingsmuligheder (kopiér, søg, analyser) \\
    \addlinespace
    F4 & AI-vejledning       & Systemet skal på baggrund af en valgt entitet generere struktureret vedligeholdelsesvejledning med trin og sikkerhedsnoter \\
    \addlinespace
    F5 & Confidence-kontrol  & Systemet skal markere resultater med lav OCR-confidence (< 60~\%) som \textit{Manuel kontrol påkrævet} \\
    \addlinespace
    F6 & Fejlhåndtering      & Systemet skal vise meningsfulde fejlbeskeder og tilbyde brugeren mulighed for at prøve igen \\
    \addlinespace
    F7 & Sprogunderstøttelse & Systemet skal understøtte vejledning på dansk og engelsk, styret af brugerindstilling \\
    \addlinespace
    F8 & Fritekst-chat       & Systemet skal understøtte fritekstbaseret kommunikation med AI, eventuelt med vedhæftet billede \\
    \bottomrule
  \end{tabularx}
\end{table}

\subsection{Ikke-funktionelle krav}
\label{subsec:ikke-funktionelle-krav}

\begin{table}[H]
  \centering
  \caption{Ikke-funktionelle krav}
  \label{tab:ikke-funktionelle-krav}
  \begin{tabularx}{\textwidth}{@{}llX@{}}
    \toprule
    \textbf{ID} & \textbf{Krav} & \textbf{Beskrivelse} \\
    \midrule
    NF1 & Svartid (billede) & OCR og entitetsformatering bør returnere svar inden for 10 sekunder \\
    \addlinespace
    NF2 & Svartid (video)   & Video-pipeline bør returnere svar inden for 30 sekunder for optagelser op til 10 sekunder \\
    \addlinespace
    NF3 & Pålidelighed      & Systemet skal håndtere netværksfejl og API-timeouts gracefully uden applikationskrasj \\
    \addlinespace
    NF4 & Brugervenlighed   & Kerneflowet fra kameraåbning til vejledning skal kunne gennemføres uden forudgående instruktion \\
    \addlinespace
    NF5 & Vedligeholdelighed & Kodebasen skal følge MVVM-arkitekturmønstret og have testdækning af alle backend-endpoints \\
    \addlinespace
    NF6 & Sikkerhed         & API-nøgler må aldrig indgå i klientkoden eller versionsstyringen \\
    \bottomrule
  \end{tabularx}
\end{table}

\section{MoSCoW-prioritering}
\label{sec:moscow}

For at sikre fokus i det iterative udviklingsforløb er kravene prioriteret efter MoSCoW-modellen~\cite{moscow}. Modellen opdeler krav i fire kategorier: \textit{Must Have}, \textit{Should Have}, \textit{Could Have} og \textit{Won't Have}.

\begin{table}[H]
  \centering
  \caption{MoSCoW-prioritering af systemkrav}
  \label{tab:moscow}
  \begin{tabularx}{\textwidth}{@{}lX@{}}
    \toprule
    \textbf{Prioritet} & \textbf{Krav} \\
    \midrule
    \textbf{Must Have}
    & Kameraoptagelse og multipart-upload (F1) \newline
      Google Cloud Vision OCR-integration (F2) \newline
      Claude Haiku entitetsformatering med JSON-output (F3) \newline
      AI-genereret vedligeholdelsesvejledning via \texttt{/api/v1/analyze} (F4) \newline
      Confidence-baseret flagging ved lav OCR-kvalitet (F5) \newline
      Fejlhåndtering med retry-mulighed (F6) \\
    \addlinespace
    \textbf{Should Have}
    & Videounderstøttelse med ffmpeg frame-udtræk (F1, F2) \newline
      Sprogvalg dansk/engelsk i brugergrænsefladen (F7) \newline
      Loading-tilstand med kontekstuelle hints under API-kald \newline
      Server health-check med statusindikator i UI \\
    \addlinespace
    \textbf{Could Have}
    & Fritekst-chat med AI via \texttt{/api/v1/chat} (F8) \newline
      Handlingsknapper per entitet (kopiér, søg, kort, analyser) \newline
      Kameraoverlay med UX-vejledning til brugeren \\
    \addlinespace
    \textbf{Won't Have}
    & SAP/ERP-integration \newline
      Android-understøttelse \newline
      Brugerautentificering og rollestyring \newline
      Offline-funktionalitet \newline
      Historik og synkronisering på tværs af enheder \\
    \bottomrule
  \end{tabularx}
\end{table}
